\documentclass[a4paper,12pt]{article}
\usepackage[utf8]{inputenc}
\usepackage[T1]{fontenc}
\usepackage[slovak]{babel}
\usepackage{geometry}
\usepackage{longtable}
\usepackage{array}
\usepackage{booktabs}

\geometry{margin=2.5cm}

\title{Užívateľské testovanie}
\date{}

\begin{document}

\maketitle

\section*{Testované časti}
Testovanie bolo zamerané na základné ovládanie aplikácie a navigáciu v nej.
\begin{itemize}
    \item Úvodná obrazovka
    \item Nastavenia
    \item Nápoveda
    \item Štatistiky
    \item Tutoriál
    \item Vloženie a export obrázka
\end{itemize}

\section*{Zameranie testovania (cieľ)}
\begin{itemize}
    \item čo ti pri tejto operácii chýbalo?
    \item čo sa stalo inak oproti očakávaniu?
    \item čo ti pri tejto operácii vadilo?
    \item čo by si zmenil?
    \item páčil sa ti dizajn?
    \item páčilo sa ti rozloženie?
\end{itemize}

\newpage
\section*{Hodnotenie použiteľnosti}
\begin{longtable}{>{\bfseries}l l}
\toprule
Hodnotenie & Popis \\
\midrule
0 & užívateľ to nenašiel vôbec \\
1 & užívateľ to našiel s veľkou nápovedou \\
2 & užívateľ to našiel s malou nápovedou \\
3 & užívateľ sa musel zamyslieť viac \\
4 & užívateľ sa musel zamyslieť \\
5 & užívateľ to našiel bez váhania \\
\bottomrule
\end{longtable}

\section*{Otázky / úlohy}

% -------------------
\subsection*{1. Úloha / Otázka}
\begin{itemize}
    \item Popis úlohy
\end{itemize}

\begin{longtable}{>{\bfseries}l c}
\toprule
Používateľ & Skóre \\
\midrule
Užívateľ 1 & \\ % doplniť
Užívateľ 2 & \\ % doplniť
Užívateľ 3 & \\ % doplniť
Výsledok & \\ % doplniť priemer
\bottomrule
\end{longtable}

\textbf{Poznámky / zistenia:}
\begin{itemize}
    \item % doplniť poznámky
\end{itemize}

\newpage
% -------------------
\subsection*{2. Úloha / Otázka}
\begin{itemize}
    \item Popis úlohy
\end{itemize}

\begin{longtable}{>{\bfseries}l c}
\toprule
Používateľ & Skóre \\
\midrule
Užívateľ 1 & \\ 
Užívateľ 2 & \\ 
Užívateľ 3 & \\ 
Výsledok & \\ 
\bottomrule
\end{longtable}

\textbf{Poznámky / zistenia:}
\begin{itemize}
    \item 
\end{itemize}

% --- Opakovať pre ďalšie otázky podľa potreby ---

\newpage
\section*{Globálne hodnotenie}
\begin{longtable}{l c}
\toprule
Úloha / Otázka & Priemer / Výsledok \\
\midrule
1. & \\ % doplniť priemer
2. & \\
3. & \\
4. & \\
5. & \\
6. & \\
7. & \\
8. & \\
\bottomrule
\end{longtable}

\textbf{Celkové priemerné hodnotenie všetkých úloh:} \\ % doplniť

\end{document}
